\chapter{Requirement Specification}

Requirement specification defines what the QuantumHop Secure Multi-Hop Communication System must perform and the conditions under which it must operate. The project is designed as a peer-to-peer encrypted messaging platform where two or more laptops on a LAN exchange quantum-secured messages using raw TCP sockets, with real-time detection of network attacks. The requirements are classified into functional requirements, non-functional requirements, hardware requirements, and software requirements.

\section{Specific Requirements}

The system must support direct message delivery between LAN peers using BB84-derived AES keys over raw TCP sockets. It must also detect and report Man-in-the-Middle, eavesdropping, and replay attacks without relying on any centralised server, cloud service, or pre-shared secret.

\subsection{Functional Requirements}

\begin{enumerate}
    \item \textbf{Peer Discovery:} The system shall automatically broadcast the machine name, IP address, and socket port over UDP to allow peers on the same LAN to discover one another without manual IP configuration.

    \item \textbf{BB84 Key Exchange Simulation:} For every outgoing message, the system shall simulate the BB84 QKD protocol by generating Alice and Bob qubit streams, performing basis selection, measuring bits, sifting matching bases, and calculating an error rate.

    \item \textbf{Shared Key Derivation:} The sifted bits from the BB84 simulation shall be hashed using SHA-256 to produce a 256-bit AES key unique to each message transmission.

    \item \textbf{AES-256-CBC Encryption:} The message payload shall be encrypted using AES-256-CBC with the BB84-derived key and a randomly generated 16-byte IV before transmission over the socket.

    \item \textbf{Raw TCP Socket Transmission:} The sender shall establish a direct TCP socket connection to the target peer's IP and port, transmit the encrypted JSON envelope, and receive a JSON response containing decryption status and error rate.

    \item \textbf{Direct Mode:} In direct mode the encrypted envelope shall be sent directly from sender to receiver with no intermediate nodes and an expected BB84 error rate near zero.

    \item \textbf{Eavesdrop Attack Mode:} In eavesdrop mode the system shall inject random basis errors into the BB84 simulation, raising the error rate above the detection threshold, causing the receiver to block the message.

    \item \textbf{MITM Relay Attack Mode:} In MITM mode the sender shall route the envelope through a relay node. The relay shall log the interception and forward the packet to the real receiver, which shall detect the relay IP and flag the message as a Man-in-the-Middle attack.

    \item \textbf{Replay Attack Mode:} In replay mode the sender shall re-transmit a previously used nonce. The receiver shall compare the nonce against a time-bound registry and reject any duplicate.

    \item \textbf{Attack Detection and Blocking:} The receiver shall block any message flagged as an attack and store it in the inbox as a blocked entry with the attack type and error rate.

    \item \textbf{Inbox Management:} All received messages, whether successfully decrypted or blocked, shall be stored in the in-memory inbox and displayed in the dashboard with sender name, message content or block reason, error rate, and timestamp.

    \item \textbf{Dashboard --- Real Network Tab:} The dashboard shall display a list of discovered network peers, a message composition panel with target peer and mode selection, and an inbox panel showing all received messages.

    \item \textbf{Dashboard --- Local Simulation Tab:} The dashboard shall provide an interactive local simulation where the user can select attack modes, run BB84 simulations, and observe the node-by-node packet flow, basis tables, and AES previews without a second machine.

    \item \textbf{Security Event Logging:} Every significant event, including BB84 completion, AES encryption, socket transmission, and attack detection, shall be logged to an in-memory event store accessible via the dashboard.

    \item \textbf{Clean Port Management:} The launcher script shall automatically terminate any processes occupying the required ports before starting the backend and frontend servers.
\end{enumerate}

\subsection{Non-Functional Requirements}

\begin{enumerate}
    \item \textbf{Security:} The system shall provide layered security through BB84 error rate monitoring, AES-256-CBC encryption, nonce-based replay prevention, and relay path detection.

    \item \textbf{Confidentiality:} Message content shall never be transmitted in plaintext. All payloads shall be AES-encrypted before socket transmission.

    \item \textbf{Integrity:} The nonce and timestamp in each envelope shall allow the receiver to detect tampered or replayed packets.

    \item \textbf{Availability:} The system shall remain operational for peer discovery and message reception continuously while the launcher script is running.

    \item \textbf{Usability:} The dashboard interface shall be operable without technical knowledge of cryptography. All attack modes shall be selectable from a dropdown, and results shall be displayed with plain-language explanations.

    \item \textbf{Performance:} BB84 simulation, AES encryption, and socket transmission shall complete within two seconds for a typical short message on a standard LAN.

    \item \textbf{Reliability:} The socket server shall handle connection errors, malformed envelopes, and handler exceptions gracefully, always returning a valid JSON error response rather than crashing silently.

    \item \textbf{Maintainability:} The backend shall be organised into separate modules: \texttt{bb84.py}, \texttt{crypto\_utils.py}, \texttt{peer\_socket.py}, \texttt{peer\_discovery.py}, \texttt{inbox.py}, \texttt{logger.py}, \texttt{attack\_detector.py}, and \texttt{config.py}.

    \item \textbf{Portability:} The system shall run on Windows, Linux, and macOS with Python 3.10 or higher and Node.js 18 or higher.

    \item \textbf{Scalability:} The peer discovery mechanism shall support any number of LAN peers, and the socket server shall handle simultaneous connections using a thread-per-connection model.
\end{enumerate}

\section{Technical Requirements}

\subsection{Hardware Requirements}

\begin{itemize}
    \item \textbf{Machines:} Minimum two laptops or desktop computers for the two-machine demonstration. One machine can be used for the local simulation mode.
    \item \textbf{Processor:} Intel Core i3 or equivalent, capable of running Python and a Node.js development server simultaneously.
    \item \textbf{RAM:} Minimum 4 GB; 8 GB recommended.
    \item \textbf{Network:} Both machines must be connected to the same Wi-Fi network, mobile hotspot, or wired LAN segment. Ports 5000 (TCP), 5010 (TCP), and 5555 (UDP) must be open in the firewall.
    \item \textbf{Storage:} Minimum 500 MB free disk space for project files and Node.js dependencies.
\end{itemize}

\subsection{Software Requirements}

\begin{itemize}
    \item \textbf{Operating System:} Windows 10/11, Ubuntu 20.04+, or macOS 12+.
    \item \textbf{Python:} Version 3.10 or higher with \texttt{pip} for dependency management.
    \item \textbf{Python Libraries:} \texttt{flask}, \texttt{flask-cors}, \texttt{cryptography} (for AES-256-CBC), \texttt{requests}.
    \item \textbf{Node.js:} Version 18 or higher with \texttt{npm} for the Vite/React frontend.
    \item \textbf{Frontend Framework:} React 18 with Vite build tool.
    \item \textbf{Browser:} Any modern browser such as Google Chrome, Microsoft Edge, or Mozilla Firefox.
    \item \textbf{Development Tools:} Visual Studio Code or any Python/JavaScript IDE, and a terminal or PowerShell for running the launcher script.
\end{itemize}
